Profil čistého disponibilního příjmu v~závislosti na hrubé mzdě dokumentuje, jak efektivní míra zatížení tvaruje reálnou mzdovou motivaci: v~ČR jsou mezní sazby pro střední příjmy vyšší než naznačuje nominální sazba daně, protože sociální pojistné vstupuje do výpočtu lineárně bez horní meze zdanitelného základu.
Pracovníci s~mzdami v~pásmu \SIrange{80}{150}{\percent} mediánu čelí nejvyšší efektivní marginální sazbě odvodu ze svého pracovního úsilí --- což oslabuje motivaci k~přesčasové práci nebo kariérnímu postupu bez skokovitého zvýšení hrubé mzdy.
Kolektivní vyjednávání může tuto past překlenout tarifními skoky, které jsou dostatečně velké na to, aby přinesly pracovníkovi čistý nárůst příjmu i~po zvýšení odvodu --- tedy zajistit, že každý tarifní stoupek do vyšší třídy přinese pozitivní motivační efekt.
Státy s~vysokým pokrytím \ac{KS} (AT, DK) vykazují plošší profil disponibilního příjmu v~nižším pásmu, protože sektorové smlouvy zajišťují pravidelné přechody mezi tarify; v~ČR tento mechanismus z~velké části chybí.

Profil čistého disponibilního příjmu v~závislosti na hrubé mzdě dokumentuje, jak efektivní míra zatížení tvaruje reálnou mzdovou motivaci: v~ČR jsou mezní sazby pro střední příjmy vyšší než naznačuje nominální sazba daně, protože sociální pojistné vstupuje do výpočtu lineárně bez horní meze zdanitelného základu.
Pracovníci s~mzdami v~pásmu \SIrange{80}{150}{\percent} mediánu čelí nejvyšší efektivní marginální sazbě odvodu ze svého pracovního úsilí --- což oslabuje motivaci k~přesčasové práci nebo kariérnímu postupu bez skokovitého zvýšení hrubé mzdy.
Kolektivní vyjednávání může tuto past překlenout tarifními skoky, které jsou dostatečně velké na to, aby přinesly pracovníkovi čistý nárůst příjmu i~po zvýšení odvodu --- tedy zajistit, že každý tarifní stoupek do vyšší třídy přinese pozitivní motivační efekt.
Státy s~vysokým pokrytím \ac{KS} (AT, DK) vykazují plošší profil disponibilního příjmu v~nižším pásmu, protože sektorové smlouvy zajišťují pravidelné přechody mezi tarify; v~ČR tento mechanismus z~velké části chybí.

Profil čistého disponibilního příjmu v~závislosti na hrubé mzdě dokumentuje, jak efektivní míra zatížení tvaruje reálnou mzdovou motivaci: v~ČR jsou mezní sazby pro střední příjmy vyšší než naznačuje nominální sazba daně, protože sociální pojistné vstupuje do výpočtu lineárně bez horní meze zdanitelného základu.
Pracovníci s~mzdami v~pásmu \SIrange{80}{150}{\percent} mediánu čelí nejvyšší efektivní marginální sazbě odvodu ze svého pracovního úsilí --- což oslabuje motivaci k~přesčasové práci nebo kariérnímu postupu bez skokovitého zvýšení hrubé mzdy.
Kolektivní vyjednávání může tuto past překlenout tarifními skoky, které jsou dostatečně velké na to, aby přinesly pracovníkovi čistý nárůst příjmu i~po zvýšení odvodu --- tedy zajistit, že každý tarifní stoupek do vyšší třídy přinese pozitivní motivační efekt.
Státy s~vysokým pokrytím \ac{KS} (AT, DK) vykazují plošší profil disponibilního příjmu v~nižším pásmu, protože sektorové smlouvy zajišťují pravidelné přechody mezi tarify; v~ČR tento mechanismus z~velké části chybí.

Profil čistého disponibilního příjmu v~závislosti na hrubé mzdě dokumentuje, jak efektivní míra zatížení tvaruje reálnou mzdovou motivaci: v~ČR jsou mezní sazby pro střední příjmy vyšší než naznačuje nominální sazba daně, protože sociální pojistné vstupuje do výpočtu lineárně bez horní meze zdanitelného základu.
Pracovníci s~mzdami v~pásmu \SIrange{80}{150}{\percent} mediánu čelí nejvyšší efektivní marginální sazbě odvodu ze svého pracovního úsilí --- což oslabuje motivaci k~přesčasové práci nebo kariérnímu postupu bez skokovitého zvýšení hrubé mzdy.
Kolektivní vyjednávání může tuto past překlenout tarifními skoky, které jsou dostatečně velké na to, aby přinesly pracovníkovi čistý nárůst příjmu i~po zvýšení odvodu --- tedy zajistit, že každý tarifní stoupek do vyšší třídy přinese pozitivní motivační efekt.
Státy s~vysokým pokrytím \ac{KS} (AT, DK) vykazují plošší profil disponibilního příjmu v~nižším pásmu, protože sektorové smlouvy zajišťují pravidelné přechody mezi tarify; v~ČR tento mechanismus z~velké části chybí.

\input{texparts/python/cz_net_income_vs_income}



